\pdfoutput=1

\documentclass[letterpaper, 10 pt, journal, twoside]{ieeetran}
\usepackage{amsmath,amsfonts}
\usepackage{algorithm}
\usepackage{algpseudocode}

\usepackage{array}
\usepackage[caption=false,font=normalsize,labelfont=sf,textfont=sf]{subfig}
\usepackage{textcomp}
% \usepackage{stfloats}
\usepackage{url}
\usepackage{verbatim}
\usepackage{graphicx}
\usepackage{cite}
\usepackage{xcolor,soul}
\usepackage{graphicx}
\usepackage{amsmath}
%\usepackage{siunitx}
%\usepackage{color}
\usepackage{cite}
\usepackage{url}

\usepackage{times}
\usepackage{fancyhdr,graphicx,amsmath,amssymb}
% \usepackage[ruled,vlined]{algorithm2e}


\usepackage{algorithm}
\usepackage{algpseudocode}
\usepackage{soul}

%\usepackage[colorlinks=true,linkcolor=blue,citecolor=blue,urlcolor=blue]{hyperref}
% ** SAVE SPACE ON FIGURES
\usepackage[font=footnotesize,belowskip=-0.5pt,aboveskip=0pt]{caption}
\usepackage{bigstrut}
\usepackage[utf8]{inputenc}
\usepackage{array}
\usepackage{multirow}
\usepackage{tabularx}
\usepackage{bigstrut}
\setlength{\extrarowheight}{5pt}
\usepackage{makecell}
\usepackage{tabu}
\usepackage{bm}
\usepackage{caption}

\usepackage{amsmath}          
\usepackage{graphics}
\usepackage{graphicx}
%\usepackage{stfloats}
\usepackage{float}
\usepackage{caption}
\usepackage{subcaption}
\usepackage{dblfloatfix}
%\usepackage{fixltx2e}
\usepackage{calc}
\usepackage{url}
\usepackage{textcomp}
% \usepackage{flushend} %comment this one out if it is messing up citations...
\usepackage{balance}
\usepackage{siunitx}
\usepackage{mdwlist}
\usepackage{comment}
\usepackage{enumerate}
% \usepackage[shortlabels]{enumitem}
\usepackage{lipsum}

% MAKE CITATIONS COMPACT
%\usepackage[nospace]{cite}
\usepackage{gensymb}

\usepackage{array}
\usepackage{multirow}
\usepackage{tabularx}
\usepackage{booktabs}

\usepackage{bigstrut}
\setlength{\extrarowheight}{5pt}
\usepackage{fixltx2e}

\usepackage{color}
\usepackage{hhline}
\usepackage[dvipsnames]{xcolor}
% \usepackage{nidanfloat }
\usepackage{float}
\usepackage{caption}
\usepackage{balance}
\usepackage[colorlinks=true, linkcolor=blue, citecolor=blue, urlcolor=blue]{hyperref}
\hyphenation{op-tical net-works semi-conduc-tor IEEE-Xplore}
\algnewcommand{\LeftComment}[1]{\Statex\(\triangleright\)~\textcolor{blue}{#1}}

% updated with editorial comments 8/9/2021

\begin{document}

\title{Beyond Point-Attached Semantics: Object-Centric Semantic Fields for Generalizable Manipulation}
% \title{A Hyperspectral Imaging Guided Robotic Grasping System with A High-Precision Hyperspectral Imaging Device}
% \title{A Hyperspectral Imaging Guided Robotic Grasping System the Custom-Designed Hyperspectral Imaging Device}
\author{Zheng Sun$^{1}$, Zhihao Li$^{1}$, Zhuo Li$^{1}$, Quentin ROUXEL $^{1}$,\\and~Fei~Chen $^{\dag 1}$,~\IEEEmembership{Senior~Member,~IEEE}
        % <-this % stops a space
%\thanks{Manuscript received: January, 23, 2025; Revised March, 24, 2025; Accepted May, 26, 2025.}
%\thanks{This paper was recommended for publication by Editor Júlia Borràs Sol upon evaluation of the Associate Editor and Reviewers' comments. This study was supported in part by the InnoHK initiative of the Innovation and Technology Commission of the Hong Kong Special Administrative Region Government via the Hong Kong Centre for Logistics Robotics.}
\thanks{
$^{1}$ Zheng Sun, Zhipeng Dong, Shixiong Wang, and Fei Chen are with the Department of Mechanical and Automation Engineering, T-Stone Robotics Institute, The Chinese University of Hong Kong, Hong Kong SAR {\tt \footnotesize zhengsun@link.cuhk.edu.hk, zhipengdong@cuhk.edu.hk,
sxwang@hklcr.hk, 
f.chen@ieee.org}.
}
\thanks{$^{2}$ Zhongyi Chu is with the School of Instrumentation and Optoelectronic Engineering, Beihang University, Beijing 100191, China  {\tt \footnotesize chuzy@buaa.edu}.}
\thanks{
$^{\dag}$ Corresponding authors
}
\thanks{Digital Object Identifier (DOI): see top of this page.}
}


\markboth{IEEE Robotics and Automation Letters. Preprint Version. May, 2025}
{SUN \MakeLowercase{\textit{et al.}}: A Hyperspectral Imaging Guided Robotic Grasping System} 
% \thanks{Manuscript received: January, 23, 2025; Revised March, 24, 2025; Accepted May, 26, 2025.}%Use only for final RAL version

\maketitle
\begin{abstract}
Generalizable robot manipulation requires stable 3D understanding of functional object parts, such as handles, tool heads, openings, and graspable regions. Raw point clouds provide geometry but lack explicit part semantics, and their sampled points vary with viewpoint, sensor configuration, and object instance. Existing 2D feature lifting and discrete 3D point-wise features enrich point clouds with semantics, but the resulting features remain attached to observation-dependent samples. We propose an object-centric continuous semantic field that conditions on an object point cloud and reads part-aware semantic embeddings at explicit 3D query locations. The field is trained from part-annotated object models and then frozen to generate semantic point clouds as object-level conditioning for manipulation policies. Experiments on RoboTwin simulation tasks and real-world bimanual object manipulation show that our representation provides more stable functional-part cues and improves policy performance over raw point-cloud, 2D feature lifting, and 3D point-wise feature baselines. Project Page: \href{https://zainzh.github.io/beyond-point-attached-semantics}{https://zainzh.github.io/beyond-point-attached-semantics}.

\end{abstract}
\begin{IEEEkeywords}
Object-Centric Representation Learning; Continuous Semantic Fields; Generalizable Robotic Manipulation
\end{IEEEkeywords}
% Two or three meaningful keywords should be added here

\section{Introduction}
\IEEEPARstart{R}{obot} manipulation is executed in 3D: grasps, contacts, and tool motions must be generated relative to object geometry observed through depth sensors. To generalize across object instances, a policy must understand functional object parts rather than object categories alone. Parts such as handles, tool heads, openings, and contact regions can vary in shape, scale, and location, while preserving their interaction roles~\cite{simeonov2022neural, pan2023tax, tang2025functo}. Learning stable part-level 3D object representations is therefore important for cross-instance generalization and for manipulation policies that can operate beyond the training object distribution.

\begin{figure*}[t]
    \vspace{-15pt}
    \centering
    \includegraphics[width=1\linewidth]{Images/teaser.png}
    \vspace{1pt}
    \caption{Teaser of our object-centric continuous semantic field. Existing 2D feature lifting and discrete 3D point-wise features attach semantics to observation-dependent samples. We instead condition a queryable semantic field on the object point cloud and read part-aware embeddings at explicit 3D query locations, producing semantic point clouds for downstream policy learning.}
    \label{fig:teaser}
\end{figure*}

Raw point clouds are a natural 3D input for manipulation policies~\cite{ze20243d, ke20243d}, but they are an unstable basis for learning such part-level cues. First, raw points encode geometry without explicit part semantics, forcing policies to infer task-relevant parts and cross-instance correspondences from demonstrations. Second, the sampled points themselves depend on viewpoint, sensor configuration, and instance geometry~\cite{goyal2023rvt}. As a result, policies trained directly on raw point clouds must learn action generation while also handling semantic ambiguity and sampling variability.

A common strategy is to enrich point clouds with semantic features. Prior works project features from 2D vision or vision-language models into 3D~\cite{chen2025g3flow, fan2026any3d, wang2024gendp}, or predict dense point-wise semantic and affordance features on the input points~\cite{chen2026learning, xu2026action}. These methods improve the semantic content of point clouds, but the features usually remain attached to the observed discrete samples. When the raw point distribution changes, the semantic representation presented to the policy changes with it. Moreover, features obtained from view-dependent images or local 3D neighborhoods can still vary with viewpoint, projection quality, and local sampling.

Our key insight is that object representations for manipulation should not merely attach semantic features to the currently observed points. Instead, they should separate object conditioning from semantic readout. We use the observed object point cloud as a geometric condition for the current instance, and use explicit 3D query positions to read out semantic embeddings. This condition-query separation allows semantic features to be generated at controlled object locations rather than only at raw sensor samples. Fig.~\ref{fig:teaser} illustrates this shift from point-attached semantic features to an object-conditioned queryable semantic field.

Based on this idea, we propose an object-centric continuous semantic field for part-aware object representation learning. The field is trained using part-annotated object models from PartNext~\cite{wang2026partnext}. In our experiments, we train one category-level field for each object family, with part labels shared across instances of that family. Given support points from an object instance, the model builds an object-conditioned tri-plane feature cache and predicts semantic embeddings and training-time part logits at explicit query locations. We supervise the queried outputs with part anchoring, cross-instance part alignment, and augmentation stability, which respectively ground predictions in part labels, align corresponding parts across instances, and improve embedding stability under support perturbations.

After training, the semantic field is frozen and used as an object-level representation module for downstream policy learning. At each policy step, we query the field at resampled object locations and pair each query coordinate with its semantic embedding to form a semantic point cloud. This representation preserves 3D spatial structure while providing functional-part cues to the policy. We use it as an additional point-cloud modality without changing the policy objective or action representation.

We evaluate our approach on four RoboTwin simulation tasks and four real-world manipulation tasks. Compared with raw point-cloud policies, 2D feature lifting, and discrete 3D point-wise features, our representation provides more stable part-level conditioning and improves policy performance, especially on tasks that require functional-part localization. Our main contributions are:
\begin{itemize}
    \item We propose an object-centric continuous semantic field that represents functional object parts as queryable 3D semantic embeddings, enabling semantic readout at explicit object locations rather than only at observed sensor points.

    \item We train category-level fields from part-annotated object models using part anchoring, cross-instance part alignment, and augmentation stability, producing embeddings that are stable and aligned across object instances.

    \item We export the learned field as policy-ready semantic point clouds and validate its benefits on multi-task and cross-instance manipulation in both RoboTwin simulation and real-world bimanual robot experiments.
\end{itemize}

\section{Related Work}

\subsection{3D Policy Learning for Robotic Manipulation}

Imitation learning policies, including diffusion-based policies, have achieved strong performance in robotic manipulation by learning action distributions from demonstrations~\cite{chi2025diffusion, zhao2023learning, florence2022implicit}. To better capture spatial interactions, recent methods condition policies on 3D observations such as point clouds~\cite{ze20243d, ke20243d}, voxel grids~\cite{shridhar2023perceiver}, or multi-view representations~\cite{goyal2023rvt, chen2025tool, li2025language, tian20253}. These methods encode raw geometry using PointNets~\cite{qi2017pointnet}, sparse 3D representations, or 3D Transformers~\cite{zhao2021point}. However, the object representation provided to the policy is often still tied to observation-dependent samples, leaving the policy to infer functional parts and cross-instance correspondences from task demonstrations. Our work focuses on this representation bottleneck: instead of changing the policy backbone, we provide a queryable, part-aware object representation as policy conditioning.

\subsection{Part-aware and Semantic Object Representations}

A large body of work enriches 3D observations with semantic or actionable features. Some methods lift features from 2D vision foundation models, such as DINOv2~\cite{oquab2023dinov2}, into 3D for open-vocabulary or language-aware perception~\cite{shen2023distilled, wang2023d, chen2025g3flow, fan2026any3d, wang2024gendp}. Other methods predict dense point-wise semantic, part, or affordance features directly on the input geometry~\cite{mo2021where2act, wu2021vat, xu2026action}. A closely related concurrent work learns part-aware 3D features for generalizable manipulation through geometric contrastive learning and semantic alignment~\cite{chen2026learning}. These methods improve the semantic expressiveness of 3D observations, but their features are primarily attached to observed discrete points.

Continuous object representations provide a queryable alternative. Neural descriptor fields and related object-centric fields have been used for correspondence, pose alignment, and category-level manipulation~\cite{simeonov2022neural, pan2023tax}. However, these fields are typically designed for descriptor matching or planning, rather than for producing part-level semantic embeddings as direct policy inputs. In contrast, our method learns an object-conditioned continuous semantic field from part-annotated object models and exports queried embeddings as semantic point clouds for downstream policy learning.

%\input{3_method}
 \section{Method}

We propose an object-centric continuous semantic field for learning part-aware object representations and using them as policy conditioning for robotic manipulation. As illustrated in Fig.~\ref{fig:method}, our method treats the observed object point cloud as a geometric condition and reads semantic embeddings at explicit 3D query locations. This separates where object geometry is observed from where semantic features are evaluated, allowing the downstream policy to condition on resampled, part-aware semantic point clouds rather than only on raw observation samples.

 \subsection{Problem Setup and Training Assumptions}

We learn a category-level object semantic field for each object family used in our experiments. Given an object point cloud, the field is conditioned on the geometry of the current object instance and can be queried at arbitrary 3D locations to produce part-aware semantic embeddings. In this work, we train the semantic field separately for each object family, while requiring the part labels to be consistent across instances within that family. The semantic field is trained using part-annotated object models from PartNext~\cite{wang2026partnext}, while the downstream policy is trained separately from task demonstrations collected in simulation or on the real robot.

For each object instance, we separate the object representation into \emph{support points} and \emph{query points}. The support set $\mathcal{P}_{sup}=\{p_i\in\mathbb{R}^3\}_{i=1}^{N_s}$ provides the geometric condition of the current object instance. During semantic field training, support points are sampled
from the object surface and used only to construct the object-conditioned field. During policy learning and deployment, support points are sampled from the observed object point cloud. The support points themselves are not the final representation consumed by the policy; they serve as the condition from which the field is built.

The query set $\mathcal{Q}=\{x_j\in\mathbb{R}^3\}_{j=1}^{N_q}$ specifies the 3D locations where semantic outputs are evaluated. During semantic field training, query points are sampled from labeled object surfaces, and each valid query point has a part label $y_j\in\{1,\ldots,C\}$ used for supervision. During policy learning and deployment, part labels are not required. We resample object locations as query points and use the field to produce semantic embeddings at these locations. Thus, support points describe the observed object instance, while query points define where semantic information is read out. This separation allows semantic features to be generated at controlled query locations rather than being restricted to the original sensor samples. Formally, the semantic field is defined as
\begin{equation}
  f_\theta(x \mid \mathcal{P}_{sup}) \rightarrow (e_x,\ell_x),
\end{equation}
where $x\in\mathbb{R}^3$ is a query coordinate, $e_x\in\mathbb{R}^d$ is an L2-normalized semantic embedding, and $\ell_x\in\mathbb{R}^C$ denotes logits over the part label set of the corresponding object category. The logits are used only for training supervision, while the semantic embedding is used as the object-level representation for downstream policy learning.


\begin{figure*}[t]
    \centering
    \includegraphics[width=1\linewidth]{Images/method.png}
   % \vspace{-10pt}
    \caption{Overview of our object-centric continuous semantic field. Support points condition an object-specific tri-plane cache; 3D query locations read out part-aware embeddings and training-time part logits. The frozen field is then queried at resampled object locations to export semantic point clouds for downstream manipulation policies.}
    %\vspace{-15pt}
    \label{fig:method}
\end{figure*}

  \subsection{Object-Conditioned Continuous Semantic Field}
  \label{sec:semantic_field}

  \paragraph{Support encoding.}
   Given the support set $\mathcal{P}_{sup}$, we extract point-wise geometric features using a frozen pre-trained 3D backbone and map them into the semantic field space with a lightweight adapter:
  \begin{equation}
      \mathcal{F}_{sup}=E(\mathcal{P}_{sup}), \qquad
      \tilde{\mathcal{F}}_{sup}=g_\phi(\mathcal{F}_{sup}),
  \end{equation}
  where $\mathcal{F}_{sup}=\{f_i\in\mathbb{R}^{d_u}\}_{i=1}^{N_s}$ and $\tilde{\mathcal{F}}_{sup}=\{\tilde{f}_i\in\mathbb{R}^{d_s}\}_{i=1}^{N_s}$. We use
  Utonia~\cite{zhang2026utonia} as $E(\cdot)$ and keep it frozen, so that semantic field learning builds on general 3D geometric priors rather than relearning low-level geometry.

  \paragraph{Tri-plane feature cache.}
  To make the representation queryable, we aggregate the adapted support features into an object-conditioned tri-plane cache:
\begin{equation}
  \mathcal{T}_{obj}
  =
  \Pi(\mathcal{P}_{sup},\tilde{\mathcal{F}}_{sup})
  =
  \{T_{xy},T_{xz},T_{yz},F_{glob}\}.
\end{equation}
Before aggregation, support coordinates are normalized to the object frame. The three planes $T_{xy},T_{xz},T_{yz}\in\mathbb{R}^{H\times W\times d_t}$ are 2D feature maps defined on the $xy$, $xz$, and $yz$ projections, respectively. The global feature $F_{glob}\in\mathbb{R}^{d_g}$ is obtained by pooling support
features. This cache is built separately for each object instance and stores the support-conditioned geometry and semantic cues used for querying.

\paragraph{Query decoding.}
Given the tri-plane cache $\mathcal{T}_{obj}$ and a query coordinate $x$, we project $x$ onto the three planes and bilinearly sample the corresponding plane features. The sampled features are fused into a local feature $F_{loc}(x)$, which is decoded together with the global feature and normalized coordinate:
\begin{equation}
  (e_x,\ell_x)
  =
  h_\theta\left(x,F_{loc}(x),F_{glob}\right).
\end{equation}
The local feature provides support-conditioned information around the query location, while $F_{glob}$ summarizes the object instance. We include the normalized coordinate $x$ as positional context in the object frame, which helps disambiguate spatially similar local patterns.


\subsection{Semantic Field Learning Objectives}
  \label{sec:method_losses}

  We supervise only the queried outputs of the field; support points are used to condition the object instance. Let $\mathcal{Q}_{valid}$ be the set of labeled
  query points and $y_x\in\{1,\ldots,C\}$ be the part label of query point $x$. For each training object, we generate two augmented support conditions $\mathcal{V}$
  and optimize a weighted sum of three objectives:
  \begin{equation}
      \mathcal{L}
      =
      \lambda_{part}\mathcal{L}_{part}
      +
      \lambda_{align}\mathcal{L}_{align}
      +
      \lambda_{stab}\mathcal{L}_{stab}.
  \end{equation}

  \paragraph{Part anchoring.}
  We implement this objective with cross-entropy supervision on the part logits:
  \begin{equation}
      \mathcal{L}_{part}
      =
      \frac{1}{|\mathcal{V}|}
      \sum_{v \in \mathcal{V}}
      \frac{1}{|\mathcal{Q}_{valid}|}
      \sum_{x \in \mathcal{Q}_{valid}}
      \mathrm{CE}(\ell_x^{(v)}, y_x).
  \end{equation}
  This objective anchors the queried outputs to the predefined part labels of the object family, making each field prediction semantically identifiable.

  \paragraph{Cross-instance part alignment.}
  Part anchoring supervises logits but does not directly organize the embedding space used by the policy. We implement this objective with supervised contrastive
  learning on the normalized embeddings. For an anchor query $i$, let $P(i)$ denote valid queries with the same part label, excluding the anchor:
  \begin{equation}
      \mathcal{L}_{align}
      =
      \frac{1}{|\mathcal{V}|}
      \sum_{v \in \mathcal{V}}
      \frac{1}{|\mathcal{I}_v|}
      \sum_{i \in \mathcal{I}_v}
      \frac{-1}{|P(i)|}
      \sum_{p \in P(i)}
      \log
      \frac{
          \exp((e_i^{(v)})^\top e_p^{(v)} / \tau)
      }{
          \sum_{a \neq i}
          \exp((e_i^{(v)})^\top e_a^{(v)} / \tau)
      },
  \end{equation}
  where $\mathcal{I}_v$ contains anchors with at least one positive sample, $P(i)$ excludes the anchor, the denominator is computed over valid query embeddings in the same batch and augmentation, and $\tau$ is a temperature parameter. This objective pulls together embeddings of the
  same part and separates different parts. Since part labels are consistent across instances within each object family, it encourages corresponding parts from
  different instances to align in the embedding space.

  \paragraph{Augmentation stability.}
  To improve stability, we enforce embedding consistency for the same query coordinate under two augmented support conditions while keeping the query coordinate fixed in the normalized object frame. Let $e_x^{(1)}$ and $e_x^{(2)}$ be the
  embeddings predicted at query point $x$ when conditioning on two augmentations of the same object instance. We minimize
  \begin{equation}
      \mathcal{L}_{stab}
      =
      \frac{1}{|\mathcal{Q}_{valid}|}
      \sum_{x \in \mathcal{Q}_{valid}}
      \left(
          1 - \langle e_x^{(1)}, e_x^{(2)} \rangle
      \right).
  \end{equation}
  This objective does not add new semantic labels; it regularizes the field to produce stable embeddings under support perturbations and small geometric transformations.
  Together, $\mathcal{L}_{part}$ anchors queried outputs to part semantics, $\mathcal{L}_{align}$ aligns corresponding parts across instances, and $\mathcal{L}_{stab}$ improves representation stability.



\subsection{Semantic Point Clouds for Policy Learning}
\label{sec:method_policy}


After semantic field training, we freeze the field and use it as an object-level representation module for policy learning. At each policy step, an observed object point cloud $\mathcal{P}_{obj}$ is used to construct the support set $\mathcal{P}_{sup}$ and the corresponding tri-plane cache. We then resample $M$ object query locations $\{x_i\}_{i=1}^{M}$ from the observed object region and evaluate the embedding branch of the field, producing a semantic point cloud
  \begin{equation}
    \mathcal{S}_{obj}=\{(x_i, f_\theta^{emb}(x_i\mid\mathcal{P}_{sup}))\}_{i=1}^{M}.
  \end{equation}

where each point stores its 3D location and queried semantic embedding. The part logits are discarded at this stage; the policy only receives semantic embeddings. For downstream manipulation, we provide these semantic point clouds as additional object-level observations. The policy observation at time $t$ is
\begin{equation}
  o_t=(\mathcal{C}_t,\{\mathcal{S}_t^k\}_{k=1}^{K},q_t),
\end{equation}
where $\mathcal{C}_t$ is the scene point cloud, $\mathcal{S}_t^k$ is the semantic point cloud of the $k$-th target object, and $q_t$ is the robot proprioceptive state. The scene point cloud preserves global geometry and spatial context, while the semantic point clouds provide part-aware object cues. We use DP3 as the
point-cloud policy backbone and treat $\mathcal{S}_t^k$ as an additional point-cloud modality. The training objective and action representation are unchanged. Thus, the policy input is generated by an explicit resampling-and-querying step, rather than being limited to features attached to the original observed points.



\begin{table}[tbp]
\centering
\small
\setlength{\tabcolsep}{3.4pt}
\caption{Simulation success rates on four RoboTwin manipulation tasks.}
\label{tab:sim_results}
\begin{tabular}{@{}lcccc@{}}
\toprule
Method & Hang Mug & Beat Hammer & Open Micro. & Put Cabinet \\
\midrule
DP3~\cite{ke20243d} & 24\% & 74\% & 22\% & 72\% \\
2D Lifting~\cite{oquab2023dinov2} & 23\% & 78\% & 21\% & 60\% \\
3D Point-wise~\cite{zhang2026utonia} & 31\% & 77\% & 26\% & 76\% \\
Ours & \textbf{37}\% & \textbf{84}\% & \textbf{35}\% & \textbf{83}\% \\
\bottomrule
\end{tabular}
\vspace{-10pt}
\end{table}

\begin{table}[tbp]
\centering
\small
\caption{Real-world success rates on four manipulation tasks.}
\label{tab:real_results}
\begin{tabular}{@{}lcccc@{}}
\toprule
Method & Grasp Mug & Beat Cube & Stir Mug & Pour Water \\
\midrule
DP3~\cite{ke20243d} & 7/20 & 8/20 & 3/20 & 4/20 \\
3D Point-wise~\cite{zhang2026utonia} & 7/20 & 9/20 & 7/20 & 5/20 \\
Ours & \textbf{17/20} & \textbf{17/20} & \textbf{10/20} & \textbf{10/20} \\
\bottomrule
\end{tabular}
\vspace{-10pt}
\end{table}

\begin{figure*}[tbp]
%\vspace{-20pt}
\centering
\includegraphics[width=0.85\linewidth]{Images/real_experiment.png}
\caption{Real-world tasks and object splits. Policies are trained on training object instances and evaluated on held-out test instances for mug grasping, tool-object contact, grasping and stirring, and pouring-related manipulation.}
\label{fig:real_world}
%\vspace{-20pt}
\end{figure*}

\section{Experiments}
We evaluate the proposed object-centric continuous semantic field through both simulation and real-world robotic experiments. Our experiments are designed to answer three questions: (1) whether the proposed representation improves the success rate of multi-task manipulation policies; (2) whether it improves policy generalization to unseen object instances; and (3) whether the learned semantic embeddings exhibit stronger cross-instance part consistency than 2D feature lifting and discrete 3D point-wise features.

\subsection{Experiment Setup}
We evaluate our method on four RoboTwin simulation tasks~\cite{chen2025robotwin} and four real-world bimanual manipulation tasks. The tasks cover mug hanging or grasping, tool-object contact, stirring, and pouring-related manipulation, all of which require localizing functional object parts such as handles, tool heads, openings, or graspable regions.

We use DP3~\cite{ke20243d} as the point-cloud imitation learning policy backbone in all experiments. In simulation, task demonstrations are generated in RoboTwin, and we compare four object representations: raw point-cloud \textbf{DP3}, \textbf{2D Feature Lifting}, \textbf{3D Point-wise Features}, and \textbf{Ours}. The DP3 baseline uses raw point clouds as 3D observations. The 2D Feature Lifting baseline projects DINOv2~\cite{oquab2023dinov2} features from RGB observations onto 3D points. The 3D Point-wise Features baseline extracts dense point-wise features from observed points using the same frozen Utonia encoder~\cite{zhang2026utonia} as our method.

For real-world experiments, demonstrations are collected on our bimanual robot using training object instances, and policies are evaluated on held-out instances with different geometry and appearance. We compare \textbf{DP3}, \textbf{3D Point-wise Features}, and \textbf{Ours}. Across all comparisons, methods share the same demonstrations, policy backbone, training pipeline, and success criteria; only the object-level visual representation differs. Detailed task definitions, object splits, hyperparameters, and success criteria are provided in the appendix.

\subsection{Quantitative Results}
\textbf{Simulation Results.}
Table~\ref{tab:sim_results} reports success rates on four RoboTwin simulation tasks. Compared with the raw point-cloud DP3 baseline, adding semantic features through 2D feature lifting or frozen 3D point-wise features does not lead to consistent improvements across tasks. For example, 2D lifting improves performance on Beat Block Hammer but degrades performance on Put Object Cabinet, while 3D point-wise features provide moderate gains but remain limited on tasks requiring precise functional-part localization. These results are consistent with our hypothesis that semantic enrichment alone is insufficient when the features remain attached to observation-dependent point samples. In contrast, our method improves performance on all four tasks, suggesting that querying an object-conditioned semantic field at resampled object locations provides more stable and policy-useful part-aware conditioning.

\textbf{Real-World Evaluation.}
Table~\ref{tab:real_results} reports real-world success rates, and Figure~\ref{fig:real_world} shows the corresponding task setup and grouped object instances. Real-world observations introduce stronger sampling variation than simulation due to sensor noise, partial views, calibration error, and appearance and geometry changes. Under these conditions, the gap between our method and the baselines becomes more pronounced. While DP3 and 3D Point-wise Features achieve limited success on several tasks, our method substantially improves performance across all four real-world tasks. This suggests that the learned semantic field provides more stable functional-part cues under real observation noise. We further analyze the cross-instance stability of the learned embeddings in Section~\ref{sec:representation_analysis}.

\begin{figure*}[t]
%\vspace{-20pt}
\centering
\includegraphics[width=1\linewidth]{Images/feature_compare.png}
%\vspace{-10pt}
\caption{Cross-instance feature visualization on real observations. Colors are obtained from a shared PCA per method and category. Compared with 2D lifting and 3D point-wise features, our field yields more consistent part-level colors across mug and hammer instances.}
%\vspace{-15pt}
\label{fig:feature_compare}
\end{figure*}

\subsection{Representation Analysis}
\label{sec:representation_analysis}
To qualitatively analyze the learned representations, Figure~\ref{fig:feature_compare} visualizes feature embeddings on four real mug instances and four real hammer instances. All features are computed from real robot observations; for clarity, we crop the target object regions for visualization. For each method and object category, we fit a shared PCA projection over embeddings from the four instances and map the first three principal components to RGB. Therefore, colors are comparable across instances within the same method-category row, but not necessarily across different methods.

The 2D Feature Lifting and 3D Point-wise Feature baselines produce meaningful local feature variations on individual instances, but the color patterns of corresponding functional parts are less consistent across instances. For example, mug handles and hammer heads do not always preserve stable colors across different object instances. In contrast, our method produces more consistent feature patterns for mug bodies and handles, as well as for hammer heads and handles. This qualitative result suggests that the learned continuous semantic field better aligns part-level embeddings across instances, providing more stable functional-part cues for downstream policy learning.

\section{Limitations}
Our method assumes rigid objects with stable geometry and semantically meaningful functional parts, and is not directly applicable to deformable objects, soft bodies, or objects whose topology changes substantially during interaction. It also relies on discrete canonical part supervision, which may be limited for ambiguous, continuous, or task-dependent functional regions. Finally, the current field mainly captures part-level semantics and does not explicitly model fine-grained geometry or physical properties. Future work could integrate geometric representations and physical attributes to support more general object understanding and manipulation.


\section{Conclusion}
We presented an object-centric continuous semantic field for learning queryable part-level object representations for manipulation. By treating the object point cloud as a geometric condition and reading semantic embeddings at explicit query positions, our method converts observation-dependent point clouds into resampleable semantic point clouds with cross-instance aligned functional-part cues. The learned field is frozen and integrated into downstream manipulation policies as an object-level conditioning module. Experiments in simulation and the real world show improved multi-task performance and cross-instance generalization over raw point-cloud and point-wise semantic baselines, suggesting that queryable object-centric semantics provide an effective representation foundation for generalizable robotic manipulation.


\section{Representation and Policy Implementation Details}
\label{app:representation_policy_details}

\paragraph{Semantic field training.}
We train the semantic field from part-annotated PartNext object models \cite{wang2026partnext}. For each object family, all training instances use a consistent part label set, such as handle/body for mugs or head/handle for hammers. The Utonia backbone is initialized from the public checkpoint and kept frozen; the trainable components are the lightweight feature adapter, tri-plane feature fusion module, and query decoder. During training, support points are sampled from the object surface to condition the field, and query points are sampled from labeled surface locations to provide part supervision. For augmentation stability, we generate two augmented support conditions from the same object using independent SO(3) rotations and Gaussian coordinate jitter, while supervising corresponding query points across the two views. Table~\ref{tab:semantic_field_hparams} summarizes the main representation training and architecture settings.

\begin{table}[h]
\centering
\small
\caption{Semantic field training and architecture settings used in our experiments.}
\label{tab:semantic_field_hparams}
\begin{tabular}{@{}p{0.32\linewidth}p{0.62\linewidth}@{}}
\toprule
Component & Setting \\
\midrule
Dataset & PartNext meshes with part annotations \\
Trainable components & Feature adapter, tri-plane fusion, and query decoder \\
Backbone & Frozen Utonia checkpoint \\
Support points & 5000 surface points per object \\
Query points & 2048 labeled surface points per object \\
Coordinate augmentation & SO(3) rotation and Gaussian jitter, $\sigma=0.005$ \\
Tri-plane cache & Resolution $64$, padding $0.05$ \\
Adapter / branch dimension & 256 / 256 \\
Semantic embedding dimension & 128, L2-normalized \\
Query decoder & Two LayerNorm--GELU MLP layers with embedding and logit heads \\
Optimizer & AdamW, learning rate $10^{-3}$, weight decay $10^{-4}$ \\
Batch size / epochs & 6 / 4000 \\
Loss weights & $\lambda_{part}=0.5$, $\lambda_{align}=0.2$, $\lambda_{stab}=0.1$ \\
Mixed precision & Enabled \\
\bottomrule
\end{tabular}
\end{table}

\paragraph{Semantic point-cloud export.}
After semantic field training, the field is frozen and used only as a representation module. For each observed object point cloud, we first remove invalid zero points and use the remaining object points as the support condition. We then sample a fixed number of object query locations and evaluate the embedding output of the field. In the default policy preprocessing, we export $256$ query points per semantic object, and each exported point stores its world-frame coordinate concatenated with a 128-dimensional semantic embedding, yielding an object-level semantic point cloud of shape $256\times(3+128)$. Part logits are not passed to the policy; they are used only during representation training and optional visualization.

\paragraph{Policy integration.}
We use DP3 as the policy backbone and keep its diffusion objective unchanged. In the observation encoder the scene point cloud and each exported semantic point cloud are treated as separate point-cloud modalities. The scene point cloud provides global geometry, while each semantic point cloud provides XYZ coordinates
concatenated with semantic embeddings for a target object. The encoded point-cloud features are concatenated with the robot proprioceptive feature to form the conditioning vector for the diffusion policy. Table~\ref{tab:policy_hparams} summarizes the downstream policy and semantic point-cloud preprocessing settings.



\begin{table}[h]
\centering
\small
\caption{Downstream policy and semantic point-cloud preprocessing settings.}
\label{tab:policy_hparams}
\begin{tabular}{@{}ll@{}}
\toprule
Component & Setting \\
\midrule
Policy backbone & DP3 point-cloud diffusion policy \\
Scene point cloud & 1024 points \\
Semantic point clouds & Up to two target objects, 256 points each \\
Semantic point dimension & $3+128$ channels \\
Observation horizon & 3 steps \\
Action horizon & 6 steps, with total horizon 8 \\
Action / proprioception dimension & 14 / 14 \\
Point-cloud encoder & PointNet encoder, output dimension 128 per modality \\
Diffusion scheduler & DDIM, 100 training steps and 10 inference steps \\
Policy optimizer & AdamW, learning rate $10^{-4}$, weight decay $10^{-6}$ \\
Batch size / epochs & 256 / 3000 \\
\bottomrule
\end{tabular}
\end{table}

\paragraph{Hardware.}
All representation and policy models are trained on a desktop workstation with a single NVIDIA RTX 4090 GPU and Ubuntu 22.04. Real-robot deployment and online policy inference are run on a laptop with an NVIDIA RTX 4070 GPU and Ubuntu 22.04. The semantic field is frozen during deployment, so online computation only requires object point-cloud extraction, field querying for the semantic point clouds, and DP3 policy inference.

\paragraph{Runtime and asynchronous execution.}
During real-world deployment, the perception and policy pipeline includes object point-cloud construction, SAM2-based object segmentation, semantic/observation encoding, and DP3 policy inference. Table~\ref{tab:runtime_breakdown} reports the measured latency of the main computation stages on the RTX 4070 laptop used for deployment. The end-to-end wall-clock time of a full inference update can be higher than the sum of these stages due to camera synchronization, robot I/O, data transfer, and action post-processing.

\begin{table}[h]
\centering
\small
\caption{Runtime breakdown of the main computation stages during real-world deployment.}
\label{tab:runtime_breakdown}
\begin{tabular}{@{}lc@{}}
\toprule
Stage & Latency \\
\midrule
Object point-cloud construction & $\sim$100 ms \\
SAM2 preprocessing and segmentation & $\sim$200 ms \\
Semantic / observation encoding & $\sim$170 ms \\
DP3 policy inference & $\sim$155 ms \\
\bottomrule
\end{tabular}
\end{table}

Since the full perception-to-policy pipeline is slower than the low-level robot command rate, we use asynchronous execution instead of blocking robot control on every policy query. A low-frequency inference thread reads RGB-D observations, extracts object point clouds, queries the semantic field, and predicts an action chunk from the policy. In parallel, a high-frequency control thread consumes the latest available action chunk and sends smooth joint commands to the robot at a fixed servo rate. When a new chunk is not yet available, the controller holds or repeats the most recent target command as a keep-alive action. We further apply command smoothing and per-step motion limits before execution. This design decouples expensive semantic perception and policy inference from low-level command streaming, enabling continuous robot motion while using the proposed semantic field in the real-world control loop.



\bibliographystyle{IEEEtran}
\bibliography{reference}
\end{document}